The estimation damage caused by statistically stealthy attacks depends on the information retained in communication records. This study examines the increase in true mean-square error under joint constraints on the modification amplitude and the total variation distance between attacked and nominal record distributions. A linear Gaussian system with full-state noisy measurements and an invertible state transition matrix is considered. Transmission events and feedback of the receiver's actual prediction are trusted, while transmitted payloads may be modified by a causal attacker. The damage is first decomposed into the mean squared conditional bias and the change in the expected nominal covariance. A covariance-weighted statistic of the next silence event then relates the conditional bias to the record distribution, yielding an upper bound for all admissible attacks. A two-step private-sign attack is constructed to obtain a matching lower bound while satisfying the complete joint-log constraint, including continuous payloads. For a fixed model, finite dimension, target time, and prescribed finite amplitude range, the worst-case damage has order $\Theta(\min\{b^2,\varepsilon\})$, where $b$ is the modification bound and $\varepsilon$ is the total variation budget. This result holds for complete logs, event histories, and exact covariance-marked records that include the next event. Counterexamples show that positive damage can coexist with unchanged record distributions when the log ends at the target time or historical information is removed. Simulations on two coupled two-dimensional systems verify the damage decomposition, examine budget-feasible attacks, and illustrate the role of the observation horizon. The bounds match in order, although their numerical constants remain conservative.
